%%%%%%%%%%%%%%%%%%%%%%%%%%%%%%%%%%%%%%%%%%%%%%%%%%%%%%%%%%%%%%%%%%%%%%%%%%%%%%%%%%%%%%%%%%%%%%%%%%%%%%%%%%%%%%%%%%%%
%%% Template para de LaTeX para o evento 12º Congresso Iberoamericano de Acústica 
%%%                      em conjunto com XXIX Encontro da Sobrac
%%% Baseado no modelo da Revista Acústica e Vibrações da Sobrac
%%% Release 04/04/2022
%%%	Desenvolvido por por Prof. William D'Andrea Fonseca, Dr. Eng. - Engenharia Acústica UFSM
%%% will.fonseca@eac.ufsm.br
%%%%%%%%%%%%%%%%%%%%%%%%%%%%%%%%%%%%%%%%%%%%%%%%%%%%%%%%%%%%%%%%%%%%%%%%%%%%%%%%%%%%%%%%%%%%%%%%%%%%%%%%%%%%%%%%%%%%
% \usepackage{carlito}
\usepackage{helvet}
%%%%%%%%%%%%%%%%%%%%%%%%%%%%%%%%%%%%%%%%%%%%%%%%%%%%%%%%%%%%%%%%%%%%%%%%%%%%%%%%%%%%%%%%%%%%%%%%%%%%%%%%%%%%%%%%%%%%
% Check if the compilation is a PDF
\usepackage{iftex}    
\ifpdf \input{glyphtounicode} \pdfminorversion=7 \pdfimageresolution=300 
\pdfgentounicode=1\fi 
%%%%%%%%%%%%%%%%%%%%%%%%%%%%%%%%%%%%%%%%%%%%%%%%%%%%%%%%%%%%%%%%%%%%%%%%%%%%%%%%%%
%%% Para documentação de códigos
\usepackage{Codes2Latex}
%%%%%%%%%%%%%%%%%%%%%%%%%%%%%%%%%%%%%%%%%%%%%%%%%%%%%%%%%%%%%%%%%%%%%%%%%%%%%%%%%%
% Para anotações
\usepackage{soul,soulutf8}
\usepackage[colorinlistoftodos,prependcaption,portuguese,textsize=scriptsize]{todonotes}
\newcommand{\hw}[1]{{\sethlcolor{lime}\hl{#1}\xspace}\addcontentsline{tdo}{todo}{#1}}
\newcommand{\hL}[1]{{\sethlcolor{yellow}\hl{#1}\xspace}\addcontentsline{tdo}{todo}{#1}}
\setlength{\marginparwidth}{2cm} 
%%%%%%%%%%%%%%%%%%%%%%%%%%%%%%%%%%%%%%%%%%%%%%%%%%%%%%%%%%%%%%%%%%%%%%%%%%%%%%%%%%
% To position the heading
\usepackage[absolute,overlay]{textpos}
  \setlength{\TPHorizModule}{1mm}
  \setlength{\TPVertModule}{1mm}
%%% Para ajuste de língua
\newcommand{\txtEN}{\selectlanguage{english}} 
\newcommand{\txtBR}{\selectlanguage{brazil}}
\newcommand{\txtSP}{\selectlanguage{spanish}}
%%%%%%%%%%%%%%%%%%%%%%%%%%%%%%%%%%%%%%%%%%%%%%%%%%%%%%%%%%%%%%%%%%%%%%%%%%%%%%%%%%
%%% Ajuste de página e margens
\usepackage[top=2cm, bottom=2cm, left=1.8cm, right=1.8cm, headsep=5mm]{geometry}
% \usepackage{showframe}
\setlength{\columnsep}{1.0cm}
\setlength{\parskip}{6pt}%
\setlength{\parindent}{0pt}%
\usepackage{setspace}
\singlespacing	% Espaço simples entre linha
%%%%%%%%%%%%%%%%%%%%%%%%%%%%%%%%%%%%%%%%%%%%%%%%%%%%%%%%%%%%%%%%%%%%%%%%%%%%%%%%%%
%%% Tipografia e hifenização
\usepackage[final,kerning=true,spacing=true]{microtype}
\microtypecontext{spacing=nonfrench}
\widowpenalty=1000 \clubpenalty=1000 \tolerance=3000

%\usepackage{etoolbox} %% Para não quebrar as referências
%\makeatletter
%\AtBeginEnvironment{thebibliography}{%
	 %\interlinepenalty=10000
   %\clubpenalty10000
   %\@clubpenalty \clubpenalty
   %\widowpenalty10000}
%\makeatother

%%% Nota de rodapé %%%%%%%%%%%%%%%%%%%%%%%%%%%%%%%%%%%%%
\usepackage[bottom,flushmargin,hang,multiple]{footmisc}

%%%%%%%%%%%%%%%%%%%%%%%%%%%%%%%%%%%%%%%%%%%%%%%%%%%%%%%%
%%% Matemática e equações
\usepackage{amsmath,amsfonts,amssymb}
\usepackage{bigints,nicefrac}
\usepackage{upgreek}
\usepackage{relsize}

% Espaço envolta das equações
\AtBeginDocument{%
 \abovedisplayskip=10pt plus 3pt minus 9pt
 \abovedisplayshortskip=5pt plus 3pt
 \belowdisplayskip=12pt plus 3pt minus 9pt
 \belowdisplayshortskip=7.5pt plus 3pt minus 4pt
}
%%%%%%%%%%%%%%%%%%%%%%%%%%%%%%%%%%%%%%%%%%%%%%%%%%%%%%%%
%%% Figuras e floats
\usepackage{graphicx,color}
\usepackage[labelformat=parens]{subfig}
\usepackage{xcolor}   
\usepackage{xspace,fancyhdr,float,calc}
 \setlength{\textfloatsep}{10pt}
 \setlength{\intextsep}{10pt}
 \setlength{\floatsep}{10pt}

\usepackage{stfloats} \fnbelowfloat %%% Para usar [b] em figure* e table*
\usepackage{multicol, wasysym, shapepar}
%%% To correct continued float %%%%%%%%%%%%%%%%%%%%%%%%%
\makeatletter
\newcommand\resetsubfigs{\setcounter{sub\@captype}{0}}
\makeatother
%\renewcommand{\figurename}{Figura:}

\usepackage{enumitem} % http://ctan.org/pkg/enumitem
%%%%%%%%%%%%%%%%%%%%%%%%%%%%%%%%%%%%%%%%%%%%%%%%%%%%%%%%
%% Links e URL
\ifpdf\usepackage{pdfpages}\fi
\PassOptionsToPackage{hyphens}{url}
\usepackage[hyperindex,pagebackref=false,pdfusetitle,pdfa=true]{hyperref}
\hypersetup{plainpages=false,colorlinks,breaklinks,linktocpage,hypertexnames=true,naturalnames=false}
\hypersetup{linkcolor=blue,citecolor=blue,filecolor=blue,urlcolor=blue} 
\urlstyle{same}
\usepackage{url, xurl}  % To break URLS
\usepackage{bookmark}
%%% References %%%%%%%%%%%%%%%%%%%%%%%%%%%%%%%%%%%%%%%%%
% Bibliografia
\usepackage{natbib} % The natbib package sorts the numbers and detects consecutive sequences, so creating “[2–4,6]”.
\setcitestyle{square,sort&compress,comma,numbers}
\usepackage{hypernat} % Para comprimir e ainda citar
\usepackage{notoccite}
\setlength{\bibsep}{4pt plus 0.3ex}
\newcommand*{\doi}[1]{doi:~\href{http://dx.doi.org/#1}{ #1}}
\newcommand{\fontrefs}{\par\fontsize{10}{12}\selectfont}
%%%%%%%%%%%%%%%%%%%%%%%%%%%%%%%%%%%%%%%%%%%%%%%%%%%%%%%%
%%% Caption
\usepackage[font=footnotesize,labelfont=bf,justification=centering, margin=4pt,parskip=0pt]{caption}
\DeclareSubrefFormat{myparens}{#1~(#2)}
\captionsetup[subfloat]{subrefformat=myparens,font=footnotesize,labelfont=bf,captionskip=10pt}
\captionsetup[table]{belowskip=0pt,aboveskip=5pt}	
\captionsetup[figure]{belowskip=0pt,aboveskip=14pt}

%%%%%%%%%%%%%%%%%%%%%%%%%%%%%%%%%%%%%%%%%%%%%%%%%%%%%%%%
%%% Configuração de títulos e seções
\usepackage{titlesec}
\titleformat{\section}
  {\normalfont\fontsize{11.5}{13.5}\bfseries}{\sffamily\thesection.}{0.5em}{\sffamily}
\titleformat{\subsection}
  {\normalfont\fontsize{10.5}{12.5}\bfseries}{\sffamily\thesubsection}{0.5em}{\sffamily}
\titleformat{\subsubsection}
  {\normalfont\fontsize{10.5}{12.5}\bfseries}{\sffamily\thesubsubsection}{0.5em}{\sffamily}
	
\titlespacing*{\section}{0pt}{6pt}{0pt}
\titlespacing*{\subsection}{0pt}{6pt}{0pt}
\titlespacing*{\subsubsection}{0pt}{6pt}{0pt}

%%% Tabelas e quadros %%%%%%%%%%%%%%%%%%%%%%%%%%%%%%%%%%%%%%%%%
\usepackage{booktabs,colortbl}
\usepackage{threeparttable}
\newcommand{\ratb}[1]{\renewcommand{\arraystretch}{#1}}
\newcolumntype{P}[1]{>{\raggedright\arraybackslash}p{#1}}
\newcolumntype{C}[1]{>{\centering\let\newline\\\arraybackslash\hspace{0pt}}m{#1}}
\newcolumntype{L}[1]{>{\raggedright\let\newline\\\arraybackslash\hspace{0pt}}m{#1}}
\newcolumntype{R}[1]{>{\raggedleft\let\newline\\\arraybackslash\hspace{0pt}}m{#1}}
\newcolumntype{V}{>{\centering\arraybackslash} m{.4\linewidth} }
\newcolumntype{A}[2]{%
    >{\minipage{\dimexpr#1\linewidth-2\tabcolsep-#2\arrayrulewidth\relax}\vspace\tabcolsep}%
    c<{\vspace\tabcolsep\endminipage}}

\newcommand{\mymulticolumn}[3]{\multicolumn{#1}{>{\columncolor{\RowColorName}}#2}{#3}}
\newcommand{\SetRowColor}[1]{\noalign{\gdef\RowColorName{#1}}\rowcolor{\RowColorName}}
%% Some colors
\definecolor{framecolor}{RGB}{226,215,255}	
% Colors for tables
\definecolor{toptb}{RGB}{221,239,255}
\definecolor{lighttb}{RGB}{255,247,221}
\definecolor{lightlinetb}{RGB}{208,239,255}
\definecolor{purple}{RGB}{112,48,160}
\definecolor{gold}{RGB}{255,192,0}
\definecolor{LightGreen}{RGB}{235,255,251}
\definecolor{LightRed}{RGB}{255,237,237}
\definecolor{LightOrange}{RGB}{255,250,237}
\definecolor{LightOrange2}{RGB}{255,235,193}
\definecolor{Green02}{rgb}{0,0.5,0}
\definecolor{Orange02}{rgb}{1,0.6,0.2}
\definecolor{LightBlue}{RGB}{212,237,255}
\definecolor{FIABlue}{RGB}{95,201,225}

%%% Matemática %%%%%%%%%%%%%%%%%%%%%%%%%%%%%%%%%%%%%%%%%%%
\DeclareMathOperator{\ji}{j}
\DeclareMathOperator{\dt}{d}
\DeclareMathOperator{\e}{e}
\DeclareMathOperator{\xp}{exp}

\renewcommand{\Re}{\mbox{Re}}
\renewcommand{\Im}{\mbox{Im}}
\newcommand{\sen}{\operatorname{sen}} 
\newcommand{\senh}{\operatorname{senh}}
\newcommand{\hada}{\operatorname{HAD}}
\newcommand{\F}{\ensuremath{\mathfrak{F}}} % F de Fourier

%%%%%%%%%%%%%% Escrever texto dentro das equacoes
\newcommand{\tx}[1]{\text{#1}}								  % Uso \tx{}
\newcommand{\txu}[1]{\ensuremath{_{\text{#1}}}} % Uso \txu{}	
\newcommand{\txup}[1]{\ensuremath{^{\text{#1}}}}% Uso \txup{}	

%%%%%%%%%%%%%% Comandos para o texto
%% Aliases
\newcommand{\dg}{{\ensuremath{^\circ}}}  

\DeclareRobustCommand{\eg}{e.g.\@\xspace}
\DeclareRobustCommand{\ie}{i.e.\@\xspace}

\makeatletter
\DeclareRobustCommand{\etc}{%
    \@ifnextchar{.}%
        {etc}%
        {etc.\@\xspace}%
}
\DeclareRobustCommand{\etal}{%
    \@ifnextchar{.}%
        {{et~al}}%
        {{et~al.}\@\xspace}%
}
\makeatother

%%%%%%%%%%%%%% Raiz quadrada com ``fechamento''
% Uso:
% \sqrt{1982} ou \oldsqrt{1982}

% New definition of square root:
% it renames \sqrt as \oldsqrt
\let\oldsqrt\sqrt
% it defines the new \sqrt in terms of the old one
\def\sqrt{\mathpalette\DHLhksqrt}
\def\DHLhksqrt#1#2{%
\setbox0=\hbox{$#1\oldsqrt{#2\,}$}\dimen0=\ht0
\advance\dimen0-0.2\ht0
\setbox2=\hbox{\vrule height\ht0 depth -\dimen0}%
{\box0\lower0.4pt\box2}}

%%%% Cria o tipo de objeto quadro %%%%%%%%%%%%%%%%%%
\usepackage{newfloat} 
\DeclareFloatingEnvironment[
    fileext=loq,
    listname={Lista de Quadros},
    name=Quadro,
    placement={!h},
    within=none,
]{quadro}

%%%%%%%%%%%%%%%%%%%%%%%%%%%%%%%%%%%%%%%%%%%%%%%%%%%%%%%%%%%%%%%%%%%%%%%%%%%%%%%%%%%%%%%%%%%%
%%% Comando para automação

%% Language commands

\providecommand{\PalavrasChaveBR}{Palavras-chave}
\providecommand{\ResumoBR}{Resumo}
\newcommand{\Resumen}[1]{\renewcommand{\ResumoBR}{Resumen} \renewcommand{\PalavrasChaveBR}{Palabras claves}}

%% Data commands

\providecommand{\NumDataRevista}{}
\newcommand{\Revista}[1]{\renewcommand{\NumDataRevista}{#1}}

\providecommand{\AutoresArtigo}{}
\newcommand{\Authors}[1]{\renewcommand{\AutoresArtigo}{#1}}

\providecommand{\AutoresFiliacoesArtigo}{}
\newcommand{\AuthorsAffiliations}[1]{\renewcommand{\AutoresFiliacoesArtigo}{#1}}

\providecommand{\FiliacoesArtigo}{}
\newcommand{\Affiliations}[1]{\renewcommand{\FiliacoesArtigo}{#1}}

\providecommand{\TituloCompletoArtigo}{}
\newcommand{\TittleComplete}[1]{\renewcommand{\TituloCompletoArtigo}{#1}}

\providecommand{\TituloCurtoArtigo}{}
\newcommand{\TitleShort}[1]{\renewcommand{\TituloCurtoArtigo}{#1}}

\providecommand{\TitleArticleArtigo}{}
\newcommand{\TitleEnglish}[1]{\renewcommand{\TitleArticleArtigo}{#1}}

\providecommand{\ResumoArtigo}{}
\newcommand{\Resumo}[1]{\renewcommand{\ResumoArtigo}{#1}}

\providecommand{\AbstractArtigo}{}
\newcommand{\Abstract}[1]{\renewcommand{\AbstractArtigo}{#1}}

\providecommand{\PalavrasChaveArtigo}{}
\newcommand{\PalavrasChave}[1]{\renewcommand{\PalavrasChaveArtigo}{#1}}

\providecommand{\KeywordsArtigo}{}
\newcommand{\Keywords}[1]{\renewcommand{\KeywordsArtigo}{#1}}

\providecommand{\PACSArtigo}{}
\newcommand{\PACS}[1]{\renewcommand{\PACSArtigo}{#1}}

% Manage complete English paper

\providecommand{\PACSEnglish}{}
\newcommand{\PACSabstract}{\renewcommand{\PACSEnglish}{
		%%%%%%%%%%%%%%%%%%%%%%%%% PACs:
		\vspace{0.7\baselineskip}
		{\fontsize{11}{11}\selectfont
		\textbf{PACS: }
		\PACSArtigo.
		\par}	
    }
}

\providecommand{\EnglishTitle}{
\sffamily\textbf{\TitleArticleArtigo}\par
\vspace{0.5\baselineskip}
}
\newcommand{\SuppressTitle}{\renewcommand{\EnglishTitle}{}}

\providecommand{\ResumoTexto}{
		{\sffamily\textbf{\ResumoBR}} \vspace{5pt}
		
		{\fontsize{11}{12.5}\selectfont 
		\ResumoArtigo
		\par}
		%%%%%%%%%%%%%%%%%%%%%%%%% Palavras-chave
		\vspace{0.7\baselineskip} \fontsize{11}{12}\selectfont
		\textbf{\PalavrasChaveBR: }{\fontsize{11}{12}\selectfont 
		\PalavrasChaveArtigo.
		\par}
		%%%%%%%%%%%%%%%%%%%%%%%%% PACs:
		\vspace{0.7\baselineskip}
		{\fontsize{11}{11}\selectfont
		\textbf{PACS: }
		\PACSArtigo.
		\par}					
    \vspace{6mm}
}
\newcommand{\SuppressResumo}{\renewcommand{\ResumoTexto}{}\SuppressTitle\PACSabstract}

%%%% PDF Info
\providecommand{\Metadata}{
\hypersetup{pdftex,
            pdfauthor={\AutoresArtigo},
            pdftitle= {\TituloCompletoArtigo},
            pdfsubject={PACS: \PACSArtigo},
            pdfkeywords={\PalavrasChaveArtigo{}.\, Keywords:  \KeywordsArtigo{}.
						\quad\quad Abstract:\,\AbstractArtigo},
            pdfproducer={LaTeX: FIA 2020 | XXIX Sobrac},
            pdfcreator={LaTeX with PDFLatex},
			pdfstartview={FitH}}
}
%%%%%%%%%%%%%%%%%%%%%%%%%%%%%%%%%%%%%%%%%%%%%%%%%%%%%%%%%%%%%%%%%%%%%%%%%%%%%%%%%%%%%%%%%%%%%%%%%%%%%%%%%%%%%%%%%%%%
%%% Edição
\Revista{FIA 2020/22 | XXIX Sobrac}
%%%%%%%%%%%%%%%%%%%%%%%%%%%%%%%%%%%%%%%%%%%%%%%%%%%%%%%%%%%%%%%%%%%%%%%%%%%%%%%%%%%%%%%%%%%%%%%%%%%%%%%%%%%%%%%%%%%%
%%% Cabeçalho
\fancypagestyle{plain}{ \fancyhf{} \fancyhead[CE,CO]{} \chead{} \setlength{\headheight}{0.7cm} \setlength{\headsep}{5mm}
\fancyhead[RE]{\fontsize{9}{10}\selectfont \NumDataRevista\; \raisebox{-0.5mm}{\includegraphics[height=6mm,page=2]{FIA-logo.pdf}}}
\fancyhead[LO]{\raisebox{-0.5mm}{\includegraphics[height=6mm,page=2]{FIA-logo.pdf}}  \fontsize{9}{10}\selectfont \NumDataRevista\; }
\fancyhead[LE,RO]{\fontsize{9}{10}\selectfont \thepage}
\fancyhead[CE,CO]{\flushleft\fontsize{8.5}{10}\selectfont \quad
%% Autores e Título
%%%%%%%%%%%%%%%%%%%%%%%%%%%%%%%%%%%%%%%%%%%%%%%%%%%%%%%%%%%%%%%%%%%%%%%%%%%%%%%%%%%%%%%%%%%%%%%%%%%%%%%%%%%%%%%%%%%%
\;\; \TituloCurtoArtigo
%%%%%%%%%%%%%%%%%%%%%%%%%%%%%%%%%%%%%%%%%%%%%%%%%%%%%%%%%%%%%%%%%%%%%%%%%%%%%%%%%%%%%%%%%%%%%%%%%%%%%%%%%%%%%%%%%%%%
}
\fancyhead[CO]{\flushright\fontsize{8.5}{10}\selectfont 
%% Autores e Título
%%%%%%%%%%%%%%%%%%%%%%%%%%%%%%%%%%%%%%%%%%%%%%%%%%%%%%%%%%%%%%%%%%%%%%%%%%%%%%%%%%%%%%%%%%%%%%%%%%%%%%%%%%%%%%%%%%%%
\TituloCurtoArtigo          \quad\; 
%%%%%%%%%%%%%%%%%%%%%%%%%%%%%%%%%%%%%%%%%%%%%%%%%%%%%%%%%%%%%%%%%%%%%%%%%%%%%%%%%%%%%%%%%%%%%%%%%%%%%%%%%%%%%%%%%%%%
}
\lfoot{} \cfoot{} \rfoot{} \renewcommand{\headrulewidth}{0.4pt}
}
\let\HeadRule\headrule
\renewcommand\headrule{\color{FIABlue}\HeadRule}
%%%%%%%%%%%%%%%%%%%%%%%%%%%%%%%%%%%%%%%%%%%%%%%%%%%%%%%%%%%%%%%%%%%%%%%%%%%%%%%%%%%%%%%%%%%%%%%%%%%%%%%%%%%%%%%%%%%%
%%%%%%%%% Cabeçalho da primeira página %%%%%%%%%%%%%%%%%%%%%%%%%%%%%%%%%%%%%%%%%%%%%%%%%%%%%
\fancypagestyle{firststyle}{ 
\fancyhf{} \fancyhead[CE,CO]{} \chead{}  \setlength{\headsep}{5mm} \renewcommand{\headrulewidth}{0pt}
\fancyhead[LE,RO]{}\fancyhead[LO,RE]{} \lfoot{} \cfoot{} \rfoot{} 
}
%%%%%%%%%%%%%%%%%%%%%%%%%%%%%%%%%%%%%%%%%%%%%%%%%%%%%%%%%%%%%%%%%%%%%%%%%%%%%%%%%%%%%%%%%%%%%%%%%%%%%%%%%%%%%%%%%%%%
%%%%%%%%%%%%%%%%%%%%%%%%%%%%%%%%%%%%%%%%%%%%%%%%%%%%%%%%%%%%%%%%%%%%%%%%%%%%%%%%%%%%%%%%%%%%%%%%%%%%%%%%%%%%%%%%%%%%
% EOF						